
\subsubsection{Calibration Inversion Metric}

For each task, a calibration inversion score is computed:

$$\text{inv}(t) = P(\text{wrong}_t) - P(\text{correct}_t)$$

Tasks with inv(t) > 0.1 exhibit calibration inversion.

\subsubsection{Calibration Clustering}

Tasks are clustered using K-means with silhouette validation. The optimal configuration is k=2 with silhouette=0.6016, revealing two behavioral clusters.

\subsubsection{Mechanism Verification Framework}

The mechanism is decomposed into five testable sub-hypotheses:

\begin{table*}[htbp]
\centering
\resizebox{\dimexpr 2\columnwidth+\columnsep\relax}{!}{%
\begin{tabular}{lllll}
\toprule
Hypothesis & Test & Threshold & Justification & Gate \\
\midrule
H-E1 & Silhouette score & > 0.3 & Standard clustering literature threshold for meaningful structure & MUST\_WORK \\
H-M1 & Confidence mean\_diff & < 0.1 & Small effect size threshold & MUST\_WORK \\
H-M2 & Rate difference & < 0.15 & Small effect size threshold (d < 0.2) & SHOULD\_WORK \\
H-M3 & Separation score & < 0.1 & Low separation indicating task-type conflation & SHOULD\_WORK \\
H-M4 & Correlation r & > 0.4 & Moderate effect size per Cohen's conventions & SHOULD\_WORK \\
\bottomrule
\end{tabular}
}
\end{table*}

\subsubsection{Bidirectional Feature Detection}

Three keyword-based features are used: user-belief-reference, context-dependent, hedged-answer. Known limitation: achieved only 2.3\% prevalence.
